\documentclass[11pt]{beamer}

\usetheme{Madrid}
\usecolortheme{default}

\usepackage[utf8]{inputenc}
\usepackage[spanish]{babel}
\usepackage{amsmath,amssymb}
\usepackage{booktabs}

\title[Identificación en modelos SVAR]{Entendiendo y resolviendo el problema de identificación en modelos VAR estructurales}
\author{}
\date{}

\begin{document}

\frame{\titlepage}

\begin{frame}{Contenido}
\tableofcontents
\end{frame}

\section{VAR reducido vs. VAR estructural}

\begin{frame}{VAR de forma reducida}
\begin{itemize}
    \item Cada variable se modela como función de:
    \begin{itemize}
        \item sus propios valores pasados,
        \item los valores pasados de las demás variables del sistema.
    \end{itemize}
    \item Ventajas:
    \begin{itemize}
        \item Se estima fácilmente por MCO (OLS).
        \item Es muy útil para pronosticar.
    \end{itemize}
    \item Problema clave: no permite separar el efecto que un choque en una variable tiene sobre las demás.
\end{itemize}
\end{frame}

\begin{frame}{VAR estructural (SVAR)}
\begin{itemize}
    \item Permite examinar relaciones \textbf{causales} entre variables.
    \item Usa teoría económica para imponer restricciones estructurales.
    \item Permite analizar el impacto de choques individuales sobre las demás variables del sistema.
\end{itemize}
\end{frame}

\section{Relación entre SVAR y VAR reducido}

\begin{frame}{Sistema estructural bivariado}
Consideremos dos variables endógenas $Y_1$ y $Y_2$, modeladas a partir de:
\begin{enumerate}
    \item sus valores rezagados un periodo,
    \item choques aleatorios propios $\epsilon_{1,t}$ y $\epsilon_{2,t}$.
\end{enumerate}

\begin{align*}
y_{1,t} &= \phi_{11} y_{1,t-1} + \phi_{12} y_{2,t-1} + b_{11}\epsilon_{1,t} + b_{12}\epsilon_{2,t}\\
y_{2,t} &= \phi_{21} y_{1,t-1} + \phi_{22} y_{2,t-1} + b_{21}\epsilon_{1,t} + b_{22}\epsilon_{2,t}
\end{align*}

Los choques $\epsilon_{1,t}, \epsilon_{2,t}$ son ruido blanco de media cero, no correlacionados serialmente y entre sí.
\end{frame}

\begin{frame}{La matriz $B$}
La matriz
\[
B = \begin{bmatrix} b_{11} & b_{12} \\ b_{21} & b_{22} \end{bmatrix}
\]
captura el impacto \textbf{estructural} de los choques $\epsilon_{1,t}$ y $\epsilon_{2,t}$ sobre $Y_1$ y $Y_2$.

\vspace{0.5cm}
\textbf{Problema:} $\epsilon_{1,t}$ y $\epsilon_{2,t}$ no son observables $\Rightarrow$ no podemos estimar $B$ directamente.
\end{frame}

\begin{frame}{De la forma estructural a la forma reducida}
Definimos los choques combinados:
\begin{align*}
u_{1,t} &= b_{11}\epsilon_{1,t} + b_{12}\epsilon_{2,t}\\
u_{2,t} &= b_{21}\epsilon_{1,t} + b_{22}\epsilon_{2,t}
\end{align*}

El sistema se convierte en un VAR de forma reducida:
\begin{align*}
y_{1,t} &= \phi_{11} y_{1,t-1} + \phi_{12} y_{2,t-1} + u_{1,t}\\
y_{2,t} &= \phi_{21} y_{1,t-1} + \phi_{22} y_{2,t-1} + u_{2,t}
\end{align*}

$\Phi$ se estima por MCO, pero los residuos $u_{1,t}, u_{2,t}$ no permiten identificar los efectos individuales de $\epsilon_{1,t}$ y $\epsilon_{2,t}$.
\end{frame}

\section{El problema de identificación}

\begin{frame}{La identidad clave}
En forma matricial: $U_t = B \epsilon_t$.

A partir de esta relación se obtiene la identidad central del SVAR:
\[
\Sigma_u = B B'
\]
donde $\Sigma_u = \mathbb{E}[u_t' u_t]$ es la matriz de covarianzas de los residuos reducidos.
\end{frame}

\begin{frame}{¿Por qué hay un problema de identificación?}
Con
\[
\Sigma_u = \begin{bmatrix} \sigma^2_{11} & \sigma^2_{12} \\ \sigma^2_{21} & \sigma^2_{22} \end{bmatrix}
= \begin{bmatrix} b_{11} & b_{12} \\ b_{21} & b_{22} \end{bmatrix}
\begin{bmatrix} b_{11} & b_{21} \\ b_{12} & b_{22} \end{bmatrix}
\]

se obtienen las ecuaciones:
\begin{align*}
\sigma^2_{11} &= b_{11}^2 + b_{12}^2\\
\sigma^2_{12} &= b_{11}b_{21} + b_{12}b_{22}\\
\sigma^2_{21} &= b_{11}b_{21} + b_{12}b_{22}\\
\sigma^2_{22} &= b_{21}^2 + b_{22}^2
\end{align*}

Como $\sigma^2_{12}=\sigma^2_{21}$, solo hay \textbf{3 ecuaciones únicas} para \textbf{4 incógnitas}: el modelo está \textbf{subidentificado}.
\end{frame}

\begin{frame}{Solución: imponer restricciones}
Para resolver el problema de subidentificación se necesitan ecuaciones adicionales, que provienen de \textbf{restricciones} basadas en teoría económica.

Formas comunes de restricción:
\begin{itemize}
    \item Sin efectos de corto plazo.
    \item Sin efectos de largo plazo.
    \item Restricciones de signo.
\end{itemize}

\vspace{0.3cm}
Ejemplo: la neutralidad monetaria implica que la política monetaria no tiene efectos acumulados de largo plazo sobre el PIB real.
\end{frame}

\begin{frame}{Ejemplos de estudios SVAR}
\small
\begin{tabular}{p{2.3cm}p{2.6cm}p{5.5cm}}
\toprule
\textbf{Autores} & \textbf{Tema} & \textbf{Restricciones} \\
\midrule
Kilian (2009) & Precios del petróleo & Sin impacto de corto plazo (1 mes) de demanda y precios sobre producción; producción sí afecta demanda en el corto plazo \\
Stock \& Watson (2001) & Choques de política monetaria & Inflación depende solo de rezagos; tasa de fondos federales no afecta desempleo contemporáneamente \\
Narayan et al. (2008) & Electricidad y PIB & Sin relación de corto plazo entre PIB real y consumo eléctrico \\
\bottomrule
\end{tabular}
\end{frame}

\section{Esquemas comunes de identificación}

\begin{frame}{Restricciones de corto plazo (Cholesky)}
\begin{itemize}
    \item Se asume que algunos choques no tienen efecto contemporáneo sobre una o más variables.
    \item Ejemplo: un choque de política monetaria no afecta la demanda agregada de inmediato $\Rightarrow b_{12}=0$.
\end{itemize}

\[
B = \begin{bmatrix} b_{11} & 0 \\ b_{21} & b_{22} \end{bmatrix}
\]

\begin{itemize}
    \item Ordenando las variables adecuadamente, $B$ es triangular inferior.
    \item Se obtiene mediante la \textbf{descomposición de Cholesky} de $\Sigma_u$.
\end{itemize}
\end{frame}

\begin{frame}{Restricciones de largo plazo (Blanchard--Quah)}
En forma matricial: $Y_t = \Phi Y_{t-1} + B\epsilon_t$.

El impacto de $\epsilon_t$ se acumula en el tiempo:
\[
\begin{array}{c|c}
\text{Periodo} & \text{Impacto de } \epsilon_t \\
\hline
T & B\epsilon_t \\
T+1 & \Phi B \epsilon_t \\
T+2 & \Phi^2 B \epsilon_t \\
\vdots & \vdots
\end{array}
\]

El efecto acumulado de largo plazo:
\[
\sum_{i=0}^{\infty}\Phi^i B \epsilon_t = (1-\Phi)^{-1} B\, \epsilon_t \equiv C\,\epsilon_t
\]
\end{frame}

\begin{frame}{Implementación de restricciones de largo plazo}
\begin{itemize}
    \item Se fijan en cero algunos elementos de la matriz $C = (1-\Phi)^{-1}B$.
    \item Procedimiento:
    \begin{enumerate}
        \item Definir $\Omega = CC' = (1-\Phi)^{-1}B\big[(1-\Phi)^{-1}B\big]'$.
        \item Obtener $C$ mediante la descomposición de Cholesky de $\Omega$.
        \item Recuperar $B = (1-\Phi)\,C$.
    \end{enumerate}
\end{itemize}
\end{frame}

\begin{frame}{Restricciones de signo}
\begin{itemize}
    \item Se asume que ciertos choques generan únicamente incrementos o disminuciones en una o más variables.
    \item Uso común: modelos de precios del petróleo (Kilian y Murphy, 2012).
\end{itemize}

\small
\begin{tabular}{lccc}
\toprule
\textbf{Choque} & \textbf{Producción} & \textbf{Act. económica} & \textbf{Precio del petróleo} \\
\midrule
Caída en oferta de petróleo & $-$ & $-$ & $+$ \\
Aumento en demanda agregada & $+$ & $+$ & $+$ \\
Aumento en demanda específica & $+$ & $-$ & $+$ \\
\bottomrule
\end{tabular}
\end{frame}

\begin{frame}{Algoritmo de restricciones de signo}
Se emplea un método iterativo de ensayo y error para hallar una matriz de rotación $Q$ tal que $B = PQ$, con $P$ la descomposición de Cholesky de $\Sigma_u$.

\begin{enumerate}
    \item Generar aleatoriamente una matriz $Q$.
    \item Calcular $P$ como la descomposición de Cholesky de $\Sigma_u$.
    \item Calcular $B$ y los impactos de los choques asociados.
    \item Verificar si los impactos cumplen las restricciones de signo:
    \begin{itemize}
        \item Si sí, conservar $Q$.
        \item Si no, descartar $Q$.
    \end{itemize}
\end{enumerate}
\end{frame}

\section{Resumen}

\begin{frame}{Resumen de los esquemas de identificación}
\small
\begin{tabular}{p{3cm}p{4.3cm}p{3.2cm}}
\toprule
\textbf{Método} & \textbf{Supuesto} & \textbf{Implementación} \\
\midrule
Efectos contemporáneos nulos & Algunos choques en $t$ no afectan variables en $t$ & $B$ triangular; Cholesky de $\Sigma_u$ \\
Efectos de largo plazo nulos & Algunos choques no tienen efecto acumulado de largo plazo & Cholesky de $\Omega=CC'$; $B=(1-\Phi)C$ \\
Restricciones de signo & Choques generan efectos de signo conocido & Búsqueda de $Q$ por ensayo y error \\
\bottomrule
\end{tabular}
\end{frame}

\begin{frame}{Conclusiones}
\begin{itemize}
    \item El VAR reducido es fácil de estimar, pero no separa los efectos causales de los choques.
    \item El SVAR permite identificar dichos efectos, pero enfrenta un problema de \textbf{subidentificación}: menos ecuaciones que incógnitas en $B$.
    \item La solución requiere \textbf{restricciones} fundamentadas en teoría económica: de corto plazo, de largo plazo, o de signo.
    \item La elección del esquema de identificación debe estar guiada por el marco teórico del fenómeno estudiado.
\end{itemize}
\end{frame}

\begin{frame}{Referencia}
Aptech Systems (2021, actualizado 2025). \emph{Understanding and Solving the Structural Vector Autoregressive Identification Problem}. \\
\url{https://www.aptech.com/blog/understanding-and-solving-the-structural-vector-autoregressive-identification-problem/}
\end{frame}

\end{document}
